\documentclass[12pt]{article}
\usepackage[utf8]{inputenc}
\usepackage[T1]{fontenc}
\usepackage[margin=1in]{geometry}
\usepackage{setspace}
\usepackage{xcolor}
\usepackage[colorlinks=true, urlcolor=blue!60!black, linkcolor=blue!60!black]{hyperref}
\setlength{\parindent}{2em}
\setlength{\parskip}{1.25em}
\renewcommand{\baselinestretch}{1.0}
\begin{document}
\begin{center}
{\LARGE \textit{Lawrence v. Texas}, 539 U.S. 558 (2003)\par}
\vspace{0.8em}
{\large Case Brief --- Four-Page Edition\par}
\end{center}
\vspace{1em}

\section{Memory Jogger}

Overruling \textit{Bowers v. Hardwick}, the Court held 6--3 that a Texas law criminalizing private, consensual same-sex intimacy violates the Due Process Clause, rejecting majoritarian moral disapproval as a legitimate state interest (\href{https://www.courtlistener.com/opinion/130160/lawrence-v-texas/}{CourtListener}).

\section{Detailed Case Facts}

On September 17, 1998, Harris County deputies responding to a falsely reported weapons disturbance entered John Geddes Lawrence's Houston apartment and found him and Tyron Garner engaged in a private, consensual sexual act. Both men were arrested, held overnight, and charged under Texas's ``Homosexual Conduct'' statute, \S~21.06(a), a Class C misdemeanor adopted in 1973 when Texas decriminalized opposite-sex sodomy but kept penalties for same-sex couples. Although rarely enforced against consenting adults in private, a conviction carried a permanent record, licensing consequences, and sex-offender registration in at least four other states.

\section{Procedural History}

Convicted before a justice of the peace, the men sought trial \textit{de novo} in county criminal court, which rejected their equal protection challenges; they pleaded \textit{nolo contendere} and were each fined \$200. A Texas Court of Appeals panel reversed, but the en banc court affirmed 7--2, treating \textit{Bowers} as controlling. After the Court of Criminal Appeals denied review, the Supreme Court granted certiorari on whether to overrule \textit{Bowers} and decided the case on June 26, 2003.

\section{Judicial Votes}

Reversed and remanded, \textbf{6--3}.
\begin{itemize}
\item \textbf{Majority:} Justice Kennedy, joined by Justices Stevens, Souter, Ginsburg, and Breyer.
\item \textbf{Concurring in the judgment:} Justice O'Connor (equal protection, without overruling \textit{Bowers}).
\item \textbf{Dissents:} Justice Scalia, joined by Chief Justice Rehnquist and Justice Thomas; Justice Thomas separately.
\end{itemize}

\section{Holding}

The Due Process Clause protects the liberty of consenting adults to engage in private, noncommercial intimate conduct in the home without criminal punishment. Moral disapproval alone is not a legitimate state interest, and \textit{Bowers} is overruled. The Court's caveats that the case did not involve ``minors,'' ``public conduct or prostitution,'' or government recognition of relationships are limiting dicta, and its references to foreign authorities were persuasive context only.

\section{Analysis of Opinions}

\textbf{Majority (Kennedy).} Tracing liberty from \textit{Meyer} and \textit{Pierce} through \textit{Griswold}, \textit{Eisenstadt}, and \textit{Roe}, Kennedy held that \textit{Bowers} misframed the interest as a ``fundamental right upon homosexuals to engage in sodomy,'' demeaning the relationship much as reducing marriage to intercourse would. Early sodomy laws targeted nonprocreative sex generally and did not single out same-sex couples until the 1970s. ``History and tradition are the starting point but not in all cases the ending point'': an ``emerging awareness,'' shown by state repeals and \textit{Dudgeon v. United Kingdom}, mattered. \textit{Casey} and \textit{Romer} had eroded \textit{Bowers}, and there was no reliance. Without calling the right fundamental, Kennedy held the statute ``furthers no legitimate state interest which can justify its intrusion into the personal and private life of the individual.''

\textbf{O'Connor (concurring in the judgment).} Applying ``a more searching form of rational basis review'' under \textit{Moreno}, \textit{Cleburne}, and \textit{Romer}, she found that punishing only same-sex sodomy targets gay persons as a class, while suggesting that preserving ``the traditional institution of marriage'' could rest on other interests.

\textbf{Scalia (dissenting).} The Court ignored \textit{Casey}'s \textit{stare decisis} standard, applied ``an unheard-of form of rational-basis review,'' and ``effectively decrees the end of all morals legislation,'' from bigamy to obscenity. It had ``taken sides in the culture war.''

\textbf{Thomas (dissenting).} He would vote to repeal the law as a legislator but found no general right of privacy in the Constitution.

\section{Examples: Future Applications}

\begin{enumerate}
\item \textbf{Same-Side (Invalidated): Fornication Bans.} A ban on private, consensual sex between unmarried adults, resting solely on religious morality, falls.
\item \textbf{Same-Side (Invalidated): Sexual Devices.} A ban on selling devices used by consenting adults in private, justified only by public morality, should fail.
\item \textbf{Opposite-Side (Upheld): Prostitution.} Kennedy expressly distinguished prostitution, and the state's interests in preventing trafficking and exploitation survive.
\item \textbf{Opposite-Side (Upheld): Minors and Coercion.} \textit{Lawrence} excluded minors and persons ``situated in relationships where consent might not easily be refused,'' such as prison inmates.
\item \textbf{Fence-Sitter: Adult Incest or Polygamy.} Challengers invoke \textit{Lawrence}'s autonomy principle; the state cites genetic risk, coercion, and ``abuse of an institution the law protects,'' and post-\textit{Dobbs} courts are less receptive.
\end{enumerate}

\section{Critique}

\textbf{Originalist and libertarian views.} Nelson Lund and John McGinnis condemned ``judicial hubris'' that replaced text and history with an amorphous autonomy principle, and Lino Graglia faulted reading Millian libertarianism into the Constitution; \textit{Dobbs} later restored \textit{Glucksberg}'s test, and Justice Thomas urged reconsidering \textit{Lawrence}. Randy Barnett, by contrast, praised a ``presumption of liberty.''

\textbf{Progressive and doctrinal views.} Laurence Tribe called \textit{Lawrence} ``the `fundamental right' that dare not speak its name,'' weaving due process and equal protection into a ``double helix'' that led to \textit{Windsor} and \textit{Obergefell}. Cass Sunstein noted the opinion's ambiguity about its standard of review.

\section{Key Quotations}

\begin{quote}
``Had those who drew and ratified the Due Process Clauses of the Fifth Amendment or the Fourteenth Amendment known the components of liberty in its manifold possibilities, they might have been more specific. They did not presume to have this insight. They knew times can blind us to certain truths and later generations can see that laws once thought necessary and proper in fact serve only to oppress. As the Constitution endures, persons in every generation can invoke its principles in their own search for greater freedom.''
\end{quote}
\noindent --- Justice Kennedy, Opinion of the Court, 539 U.S. at 578--79.

\begin{quote}
``Liberty protects the person from unwarranted government intrusions into a dwelling or other private places. In our tradition the State is not omnipresent in the home. And there are other spheres of our lives and existence, outside the home, where the State should not be a dominant presence. Freedom extends beyond spatial bounds. Liberty presumes an autonomy of self that includes freedom of thought, belief, expression, and certain intimate conduct. The instant case involves liberty of the person both in its spatial and more transcendent dimensions.''
\end{quote}
\noindent --- Justice Kennedy, Opinion of the Court, 539 U.S. at 562.

\begin{quote}
``Moral disapproval of this group, like a bare desire to harm the group, is an interest that is insufficient to satisfy rational basis review under the Equal Protection Clause.''
\end{quote}
\noindent --- Justice O'Connor, Concurring in the Judgment, 539 U.S. at 582.

\end{document}
